% Development addendum only. Not a replacement for the current manuscript.
% Requires amsmath, amssymb, bm. No experimental superiority is asserted here.

\subsection{Head-Side Fusion of Separately Recovered Updates}

The linear update in the current manuscript adds the recovered FL and FU
aggregates directly. The receive coefficients compensate for the channel
and control the aggregation noise; they are not task-combination weights
under exact channel alignment. We therefore evaluate an additional
head-side fusion rule based on two-task gradient surgery
\cite{Yu_NeurIPS_2020_PCGrad}.

Let $\widehat{\mathbf L}_k$ and $\widehat{\mathbf U}_k$ collect all
recovered FL and FU blocks, respectively. These vectors include both the
visual and textual adapter coordinates. For two vectors $\mathbf L$
and $\mathbf U$, define
\begin{equation}
q(\mathbf L,\mathbf U)=[-\mathbf L^{\mathsf T}\mathbf U]_+.
\end{equation}
The two-task PCGrad fusion is
\begin{equation}
\begin{aligned}
\Phi_{\rm PC}(\mathbf L,\mathbf U)
={}&\mathbf L+\mathbf U\\
&+q(\mathbf L,\mathbf U)
\left(\frac{\mathbf L}{\|\mathbf L\|_2^2}
+\frac{\mathbf U}{\|\mathbf U\|_2^2}\right).
\end{aligned}
\end{equation}
If either task vector is zero, the correction is zero. Thus, the rule
reduces to the original sum when the directions do not conflict or only
one task is active. The head applies
\begin{equation}
\boldsymbol\theta_{k+1}=\boldsymbol\theta_k+
\Phi_{\rm PC}(\widehat{\mathbf L}_k,\widehat{\mathbf U}_k).
\end{equation}
The separate estimates are necessary to evaluate their inner product
and norms. A single combined aggregate does not generally identify
these quantities. A fair combined-stream comparator applies the same
two-task surgery independently at each UAV before transmission. Local
surgery and surgery after global aggregation need not agree.

This is a generic two-objective rule applied to joint visual--textual
adapter updates. It is not a new VLM-specific gradient-surgery theorem,
nor does the nonnegative alignment of the resulting direction imply
successful unlearning or proximity to a retrained model.

\paragraph{Noise-regularized variant.}
Let $v_m^a=\mathbb E\|\mathbf e_m^a\|_2^2$ and
$s_m^a=v_m^a/d_m$ be the designed block MSE and per-coordinate noise
variance. Define $A_m=[\|\widehat{\mathbf L}_m\|_2^2-v_m^{\mathrm L}]_+$
and $B_m=[\|\widehat{\mathbf U}_m\|_2^2-v_m^{\mathrm U}]_+$, and use
\begin{equation}
\widehat\sigma_{\rm dot}^2=
\sum_m\left(s_m^{\mathrm U}A_m+s_m^{\mathrm L}B_m
+d_ms_m^{\mathrm L}s_m^{\mathrm U}\right).
\end{equation}
The registered variant replaces $q$ by
$[-\widehat{\mathbf L}^{\mathsf T}\widehat{\mathbf U}
-1.645\widehat\sigma_{\rm dot}]_+$ and replaces the two denominators
by $\|\widehat{\mathbf L}\|_2^2+\sum_m v_m^{\mathrm L}$ and
$\|\widehat{\mathbf U}\|_2^2+\sum_m v_m^{\mathrm U}$, respectively.
This is a heuristic regularization, not a certified confidence bound.
The head uses only received vectors and designed noise statistics.
Evaluation data and retrained-model predictions are not control inputs.

\paragraph{Effect on the distortion analysis.}
For a selected fusion map $\Phi_k$, the ideal update must also be changed
to $\overline{\boldsymbol\theta}_{k+1}=\boldsymbol\theta_k+
\Phi_k(\mathbf L_k,\mathbf U_k)$, where the two arguments are the
error-free global aggregates. The relevant output error is
\begin{equation}
\mathbf e_{\Phi,k}=\Phi_k(\widehat{\mathbf L}_k,
\widehat{\mathbf U}_k)-\Phi_k(\mathbf L_k,\mathbf U_k).
\end{equation}
Even with zero-mean channel noise, this error need not have zero mean.
Consequently, the existing curvature-weighted sum of component MSEs is
only a pre-fusion transport constraint; its original zero-mean
linear-update proof cannot be reused as a post-fusion quality bound.
If a task objective is $L_J$-smooth along the relevant segment and
$\|\nabla J(\overline{\boldsymbol\theta}_{k+1})\|_2\leq G_J$, a valid
generic inequality is
\begin{equation}
\begin{aligned}
\mathbb E[J(\boldsymbol\theta_{k+1})]
-J(\overline{\boldsymbol\theta}_{k+1})
\leq G_J\mathbb E\|\mathbf e_{\Phi,k}\|_2
+\frac{L_J}{2}\mathbb E\|\mathbf e_{\Phi,k}\|_2^2.
\end{aligned}
\end{equation}
A certified relation between this post-fusion error and component MSEs
requires additional analysis. The development experiment logs realized
post-fusion errors and does not claim such a certificate. Similarly,
optimizing communication subject to the old transport constraints is not
an SOCP proof for a newly imposed nonlinear post-fusion constraint.

\paragraph{Delay and energy.}
Head fusion occurs after uplink reception and before the aggregate is
broadcast. Its computation duration and energy must be included, along
with the additional hovering time. They cannot be hidden inside the
parallel local-computation interval. The development report explicitly
adds these costs using the same processor model as the original study.

% Other manuscript corrections relevant to this change:
% C6, not C7, is exact alignment in the currently pasted constraint list.
% C7 is component MSE; C8 is the pre-fusion curvature-weighted distortion.
% Do not replace the old tables with results from a different learning
% objective or transport rule without rerunning the stated comparison.
